\documentclass[10pt,a4paper]{amsart}
\usepackage[margin=19mm]{geometry}
\usepackage{amsmath,amssymb,mathtools}
\usepackage[scr=rsfs,cal=euler]{mathalfa}
\usepackage{xfrac}
\usepackage[unicode,hidelinks,bookmarks=false]{hyperref}
\newcommand{\ie}{i.e.\ }
\newcommand{\cf}{cf.\ }
\newcommand{\resp}{respectively\ }
\newcommand{\Subsection}{\subsection}
\providecommand{\nobreakdash}{\nobreak\hbox{-}}
\setlength{\emergencystretch}{2em}
\multlinegap0pt
\usepackage{cleveref}
\newtheorem{theorem}{Theorem}
\newtheorem{prop}[theorem]{Proposition}
\newtheorem{definition}[theorem]{Definition}
\newcommand{\rlasso}{\texttt{rLasso}}
\newcommand{\omp}{\texttt{OMP}}
\title[Theorem 2.8]{High-dimensional sparse recovery from function samples -- Decoders, guarantees and instance optimality (Theorem 2.8)}
\author{Moritz Moeller, Sebastian Neumayer, Kateryna Pozharska, Tizian Sommerfeld, Tino Ullrich}
\date{Source: arXiv:2503.16209v2 (CC BY 4.0)}
\begin{document}
\maketitle

\noindent\emph{Source and licence.} High-dimensional sparse recovery from function samples -- Decoders, guarantees and instance optimality, by Moritz Moeller, Sebastian Neumayer, Kateryna Pozharska, Tizian Sommerfeld, Tino Ullrich. Version arXiv:2503.16209v2
(CC BY 4.0), distributed by arXiv under the Creative Commons Attribution 4.0 International licence,
http://creativecommons.org/licenses/by/4.0/, as recorded in the arXiv licence metadata for this exact
version and shown on the abstract page https://arxiv.org/abs/2503.16209v2. Full licence notice:
\texttt{licenses/arxiv-2503.16209-CC-BY-4.0.txt}.\par\smallskip

\noindent\textit{Editorial scope.} Counted unit: the article's Theorem 2.8 (source0.tex lines 396--403, printed as Theorem 2.8): there are universal constants $\alpha,C_B>0$ such that for every $n$ and finite $J\subset I$ with $n<\lvert J\rvert$, if at least as many samples as required by the displayed sampling bound are drawn, then with probability at least $1-\gamma$ the decoder $H_{\mathbf{X}}$ satisfies the displayed $L_q$ error estimate for all $2\le q\le\infty$ and all $f$ in the weighted space. The article's complete original proof of it is delivered with it (source0.tex lines 405--442, one \texttt{proof} environment of 38 lines, adjacent to the statement, with no gap and no deferral). No mathematics is written, repaired or re-derived: the statement and the proof are the article's own lines, in the article's own notation, and no retained line is substituted, re-flowed or omitted, so nothing is removed and the package carries no omission record. Retained uncounted as the source-local context the proof needs: the paragraph defining the measurement matrix and the sample vector (lines 256--262), the definition of the square-root Lasso decoder with its minimization problem (lines 264--274, printed as Definition 2.1, source label \texttt{def:square\_root\_lasso}), the robust recovery guarantee for that decoder (lines 322--328, printed as Proposition 2.5, source label \texttt{prop:rLasso\_for\_NSP}) and the vector-level recovery error proposition (lines 371--383, printed as Proposition 2.7, source label \texttt{prop:Recovery\_vector\_errors}). Every cross-reference of every retained line resolves inside this document, and no retained line contains a citation, so the wrapper carries no bibliography and no bibliography entry. No retained line contains a figure environment, an image-inclusion command or drawing code, so no image is delivered and the package declares no asset dependency.

Editorial formatting only, in addition: an AMS \texttt{amsart} preamble that restates the article's own declarations and macros used by the retained lines -- the environments \texttt{theorem}, \texttt{prop} and \texttt{definition} on one shared counter exactly as in the article but without the article's section-based prefix; the two macros \texttt{\textbackslash rlasso} and \texttt{\textbackslash omp}, copied verbatim from the article's own preamble (source0.tex lines 51 and 52); and the standard \texttt{cleveref} package that the article itself uses for its \texttt{\textbackslash Cref} cross-references. The retained environments therefore print in this document as Definition 1, Proposition 2, Proposition 3 and Theorem 4 in order of appearance, whereas the article prints the same four items as Definition 2.1, Proposition 2.5, Proposition 2.7 and Theorem 2.8; the numbered displays of the retained lines print consecutively in order of appearance, whereas the article numbers its equations per section.

The complete native source of the article as distributed in the arXiv e-print for this exact version is bundled unchanged as \texttt{sources/175107/source0.tex}: it is the mechanical concatenation of the article's five own source files (main.tex and the four section files) in the order in which the main file reads them, byte identical to those files apart from the single header comment line the pipeline prepends. The article was typeset with the standard \texttt{article} class; no publisher class or style file is shipped in the e-print and none is redistributed or used.
We consider the two nonlinear decoders \emph{Square Root Lasso} (\rlasso) and \emph{Orthogonal Matching Pursuit} (\omp).
For the function recovery based on a dictionary $\mathcal{B}$, we fix a finite index set $J\subset I$, and define the matrix
\begin{gather}\label{eq:matrix}
	\mathbf{A} \coloneqq m^{-\sfrac{1}{2}}\left(b_j(x^l)\right)_{1 \leq l \leq m, j \in J} \in \mathbb{C}^{m\times \lvert J\rvert}
\end{gather}
and the vector $\mathbf{y} \coloneqq m^{-\sfrac{1}{2}} f(\mathbf{X}) \in \mathbb{C}^m$.
First, we specify \rlasso.

\begin{definition}[\rlasso]\label{def:square_root_lasso}
For $\lambda>0$, we define $R_{\mathbf{X}}^\lambda\colon C(\Omega) \to V_J$ by
\begin{gather}\label{eq:Lasso_Decoder}
	R_{\mathbf{X}}^\lambda(f) = \sum_{j \in J} r_j(\mathbf{y}) b_j,
\end{gather}
where $\mathbf{r}(\mathbf{y}) = [R_{\mathbf{X}}^\lambda(f)] \in \mathbb{C}^{\lvert J\rvert}$ is a (fixed) solution of the \rlasso \ minimization problem
\begin{gather}\label{eq:MinimizationProblem}
	\min_{\mathbf{z} \in \mathbb{C}^{\lvert J\rvert}} \lVert\mathbf{z}\rVert_{\ell_1} + \lambda \lVert\mathbf{A} \mathbf{z} - \mathbf{y}\rVert_{\ell_2}
\end{gather}
with the matrix $\mathbf{A}$ from \eqref{eq:matrix}, where we implicitly assume that there is a deterministic selection rule in the case of multiple solutions.
\end{definition}

\begin{prop}\label{prop:rLasso_for_NSP} Let $\mathbf{A} \in \mathbb{C}^{m\times N}$ satisfy $\mathbf{A}\in\operatorname{NSP}(n,2,\lVert\cdot\rVert_{\ell_2},\varrho,\tau)$.
Then for all $\mathbf{y} \in \mathbb{C}^m$ and $\mathbf{v} \in \mathbb{C}^{N}$, any solution $\mathbf{r}(\mathbf{y}) \in \mathbb{C}^{N}$ of the \rlasso \ minimization problem \eqref{eq:MinimizationProblem} with $\lambda = 3\tau n^{\sfrac{1}{2}}$ satisfies
\begin{equation}\label{eq:rLasso_NSP_lq}
	\lVert\mathbf{v}- \mathbf{r}(\mathbf{y}) \rVert_{\ell_q} \leq \beta\bigl(n^{\sfrac{1}{q}-1}\sigma_n(\mathbf{v})_{\ell_1} + n^{\sfrac{1}{q}-\sfrac{1}{2}}  \lVert\mathbf{A}\mathbf{v}-\mathbf{y}\rVert_{\ell_2}\bigr)
\end{equation}
for all $q\in[1,2]$ and a constant $\beta>0$ that depends only on $\varrho$ and $\tau$.
\end{prop}

\begin{prop}\label{prop:Recovery_vector_errors}
Let $\mathcal{B} \subset L_2(\mu)$ be a BOS and \smash{$B \coloneqq \sup_{j \in I}\lVert b_j\rVert_{L_\infty}$}.
For $m\in\mathbb{N}$ and $J\subset I$ finite, let the (random) matrix $\mathbf{A}\in\mathbb{C}^{m\times\lvert J\rvert}$ be defined through \eqref{eq:matrix} with \smash{$\mathbf{X}\overset{iid}{\sim} \mu$}.
There exist universal constants $\alpha, \beta>0$ such that for all $n \in \mathbb{N}$, $\mathbf{v} \in \mathbb{C}^{\lvert J\rvert}$, $\gamma \in (0,1)$ and ${\mathbf{y}} \in \mathbb{C}^m$ with 
\begin{equation}\label{eq:NumberSamples}
	m \geq \alpha B^2 n \bigl({\log^2} \bigl(B^2n\bigr) \log\lvert J\rvert + \log\bigl(\gamma^{-1}\log\bigl(B^2n\bigr)\bigr)\bigr)
\end{equation}
the \rlasso \ (with $\lambda(n)$ according to \Cref{prop:rLasso_for_NSP}) or \omp \ (after sufficiently many iterations) reconstruction vector $\mathbf{h}(\mathbf{y})\in\mathbb{C}^{\lvert J\rvert}$ satisfies
\begin{equation}\label{eq:Recovery_lq_error}
    \lVert \mathbf{v} - \mathbf{h}(\mathbf{y}) \rVert_{\ell_q} \leq \beta\bigl(n^{\sfrac{1}{q}-1}\sigma_n(\mathbf{v})_{\ell_1} + n^{\sfrac{1}{q}-\sfrac{1}{2}}  \lVert\mathbf{A}\mathbf{v}-\mathbf{y}\rVert_{\ell_2}\bigr)
\end{equation}
for all $q\in[1,2]$ with probability at least $1 - \gamma$.
\end{prop}

\begin{theorem}\label{thm:Error_Recovery_Ops}
There exists a universal constants $\alpha, C_B>0$ such that for $n \in \mathbb{N}$ and $J\subset I$ finite with $n<\lvert J\rvert$ the following holds.
If we at least as many sample points \smash{$\mathbf{X}=\{x^1,\ldots,x^m\} \overset{iid}{\sim} \mu$} as required in \eqref{eq:NumberSamples}, then with probability at least $1-\gamma$ we have
\begin{gather}\label{eq:ErrorEstimateRIP}
	\lVert f - H_{\mathbf{X}}(f)\rVert_{L_q} \leq C_B n^{\sfrac{1}{2}-\sfrac{1}{q}}\big(n^{-\sfrac{1}{2}}\sigma_n(f;\mathcal{B})_{\mathcal{A}_1} + \lVert f-P_Jf\rVert_{L_\infty}\big)
\end{gather}
for any $2\leq q\leq \infty$ and $f \in \mathcal{A}_1$.
\end{theorem}

\begin{proof}
First, we consider $q=\infty$.
Applying \eqref{eq:Recovery_lq_error} from \Cref{prop:Recovery_vector_errors} with $\mathbf{v} = P_Jf$ and $\mathbf{y} = m^{-\sfrac{1}{2}}f(\mathbf{X})$ yields that
\begin{align}\label{eq:Est_PN_term}
\begin{split}
    \lVert P_Jf - H_{\mathbf{X}}(f) \rVert_{L_\infty} &\leq \sum_{j\in J} \lVert b_j\rVert_{L_\infty} \big\lvert[P_Jf]_j - [H_{\mathbf{X}}(f)]_j \big\rvert \leq B \big\lVert[P_Jf] - [H_{\mathbf{X}} (f)]\big\rVert_{\ell_1}\\
    &\leq B\beta\big(\sigma_n(P_Jf)_{\ell_1} + n^{\sfrac{1}{2}} m^{-\sfrac{1}{2}} \lVert P_Jf(\mathbf{X}) - f(\mathbf{X})\rVert_{\ell_2}\big)\\
    &\leq B\beta\big(\sigma_n(P_Jf)_{\ell_1} + n^{\sfrac{1}{2}} \lVert P_Jf - f\rVert_{L_\infty}\big)
\end{split}
\end{align}
holds with the stated probability. Using \eqref{eq:Est_PN_term} and $ \sigma_n(P_Jf)_{\ell_1} \leq \sigma_n(f)_{\ell_1} = \sigma_n(f)_{\mathcal{A}_1}$, we further obtain
\begin{align}
\begin{split}
    \lVert f - H_{\mathbf{X}}(f)\rVert_{L_\infty} & \leq \lVert f - P_Jf \rVert_{L_\infty} + \lVert P_Jf - H_{\mathbf{X}}(f) \rVert_{L_\infty}\\
    &\leq B \beta \sigma_n(P_Jf)_{\ell_1} + (B\beta n^{\sfrac{1}{2}}+1)\lVert f - P_Jf\rVert_{L_\infty}\\
    &\leq C_{\infty} n^{\sfrac{1}{2}} \bigl(n^{-\sfrac{1}{2}}\sigma_n(f)_{\mathcal{A}_1} + \lVert f-P_Jf\rVert_{L_\infty} \bigr).
\end{split}
\end{align}
This concludes the case $q =\infty$.

For $q = 2$, we note that $\lVert\cdot\rVert_{L_2}\leq \lVert\cdot\rVert_{L_\infty}$.
Thus, the Parseval identity and \eqref{eq:Recovery_lq_error} from \Cref{prop:Recovery_vector_errors} imply
\begin{align}
\begin{split}
    \lVert f - H_{\mathbf{X}}(f)\rVert_{L_2} & \leq \lVert P_Jf -  H_{\mathbf{X}}(f)\rVert_{L_2} + \lVert f -P_Jf \rVert_{L_\infty}\\
    & = \bigl\lVert [P_Jf] - [H_{\mathbf{X}}(f)]\bigr\rVert_{\ell_2} + \lVert f - P_Jf \rVert_{L_\infty}\\
    &\leq \beta n^{-\sfrac{1}{2}}\sigma_n(P_Jf)_{\ell_1} + (\beta+1) \lVert f - P_Jf\rVert_{L_\infty}\\
    & \leq C_2 \big( n^{-\sfrac{1}{2}} \sigma_n(f)_{\mathcal{A}_1} + \lVert f-P_Jf\rVert_{L_\infty} \big).
\end{split}
\end{align}

For $2<q<\infty$, H\"older's inequality implies
\begin{gather}
\lVert f - H_{\mathbf{X}}(f)\rVert_{L_q} \leq \lVert f - H_{\mathbf{X}}(f)\rVert_{L_2}^{1-\theta}\lVert f - H_{\mathbf{X}}(f)\rVert_{L_\infty}^\theta,
\end{gather}
where $\theta$ has to be chosen such that $\sfrac{1}{q}= \sfrac{(1-\theta)}{2}+ \sfrac{\theta}{\infty}$, which yields $\theta = 1 - \sfrac{2}{q}$.
Now, the desired estimate holds with probability at least $1-\gamma$.
\end{proof}

\end{document}
